\section{Conclusion}
\label{sec:conclusion}

This paper studies adaptive weight allocation for distributed KLA-based LMB fusion under heterogeneous packet loss. The revised evidence supports a narrower conclusion than a monolithic ``all modules help'' claim. Realized link quality is the dominant factor in the tested communication-constrained scenario. Covariance quality provides an independent posterior-concentration signal and gives a further cardinality benefit when combined with link quality. The combined branch-aware refinement then supplies a small but repeatable improvement in OSPA consensus error and matched localization disagreement.

Existence confidence has little isolated incremental effect on the covariance-link backbone in the present experiment. Branch decoupling and weak structure-aware modulation are therefore evaluated together as one branch-aware refinement rather than assigned separate empirical claims. The detailed ablation makes this distinction explicit and separates architectural motivation from measured effect size.

Relative to Fixed Metropolis, the Balanced configuration reduces OSPA consensus error, matched localization disagreement, and cardinality dispersion by $28.8\%$, $35.5\%$, and $85.4\%$, respectively, with approximately $1.10\times$ fixed filtering/fusion runtime. Its truth-referenced local E-OSPA and cardinality error remain substantially below the fixed baseline, while local RMSE stays close to the covariance-link backbone.

The Cardinality-critical mode retains the same branch-aware backbone and adds FID-FIA only to the existence branch. Relative to Balanced, it reduces cardinality dispersion by $34.8\%$, local E-OSPA by $2.0\%$, and local cardinality error by $26.2\%$. The tradeoff is explicit: OSPA consensus error increases by $1.0\%$, localization disagreement by $8.9\%$, and local RMSE by $6.6\%$. Balanced is therefore the spatial-consensus and lower-cost mode, whereas Cardinality-critical is appropriate when target-number errors dominate.

The additional payload-semantics stress study clarifies the role of light and heavy communication. In the legacy moment-matched KLA path, full heavy GM-LMB messages do not change the fused result relative to light moment-matched summaries in the tested crossing scenario. After mixture-aware label-wise GM-KLA is enabled, selected heavy mixture components are preserved and close-crossing performance improves: crossing-window E-OSPA decreases from $3.298$ to $3.151$, the absolute crossing p90 and worst-case values improve, and the mixture-aware heavy arm wins $16/20$ paired seeds. This supports a bounded but important design lesson: payload hierarchy and the receiving fusion operator must be co-designed.

The current operating modes are confirmed in the main tiered-loss scenario, and the mixture-aware heavy extension is confirmed in a small high-clutter crossing stress. Repeating ideal-communication, communication-level, and arithmetic-average studies is the next validation step for the weighting method. Further work should also consider richer sensor-parameter heterogeneity, asynchronous or multi-rate communication, stronger mode-persistence management for heavy mixtures, and target-wise information measures once reliable label association across neighboring LMB posteriors is available.
